\documentclass[12pt,a4paper]{article}

\usepackage[utf8]{inputenc}
\usepackage[T1]{fontenc}
\usepackage[polish]{babel}
\usepackage[a4paper,margin=2.5cm]{geometry}
\usepackage{enumitem}
\usepackage{hyperref}
\usepackage{titlesec}
\usepackage{graphicx}
\usepackage{subcaption}
\usepackage{tabularx}
\usepackage{booktabs}
\usepackage{url}
\usepackage{comment}
\usepackage[]{hypcap}
\usepackage{chngcntr}
\usepackage{float}
\usepackage{ragged2e}

\newcommand{\projecttitle}{Modyfikacja algorytmu CMA-ES}
\newcommand{\authorone}{Aleksandra Wieczorek}
\newcommand{\authortwo}{Wojciech Wójcik}
\newcommand{\projectdate}{Warszawa, \today}


\begin{document}
\sloppy

\begin{titlepage}
    \centering

    \includegraphics[width=0.8\textwidth]{politechnika.png}\par
    \vspace{2.5cm}

    {\Large Wstęp do algorytmów ewolucyjnych\par}
    \vspace{1.2cm}

    {\Huge \bfseries \projecttitle \par}
    \vspace{2cm}

    {\large \projectdate\par}

    \vfill
    {\large \authorone\par}
    \vspace{0.3cm}
    {\large \authortwo\par}

    
\end{titlepage}
\newpage

%\section{Wprowadzenie}
%W algorytmie CMA-ES wektor p sigma , reprezentujący ścieżkę ewolucyjną, standardowo inicjalizowany jest wektorem zerowym. Zmodyfikuj algorytm tak, aby wektor ten był inicjalizowany wektorem losowym.
%Zbadaj działanie wersji standardowej algorytmu oraz wersji zmodyfikowanej, używając wybranych funkcji testowych. W badaniu wykorzystaj co najmniej dwa róże generatory liczb losowych.



\section{Opis problemu}
 
Algorytm CMA-ES (Covariance Matrix Adaptation Evolution Strategy) jest jedną z~najskuteczniejszych metod optymalizacji ciągłej bez użycia gradientu. Algorytm iteracyjnie próbkuje populację kandydatów z~wielowymiarowego rozkładu normalnego $\mathcal{N}(\mathbf{m}, \sigma^2 \mathbf{C})$, a~następnie adaptuje parametry tego rozkładu --- średnią~$\mathbf{m}$, krok mutacji~$\sigma$ oraz macierz kowariancji~$\mathbf{C}$ --- na podstawie jakości uzyskanych rozwiązań.
 
Kluczowym mechanizmem kontroli kroku $\sigma$ jest ścieżka ewolucyjna $\mathbf{p}_\sigma$ (ang.\ \textit{cumulative step-size adaptation path}). Wektor ten akumuluje informację o~kierunkach kolejnych kroków algorytmu i~jest aktualizowany w~każdej generacji zgodnie ze wzorem:
\[
\mathbf{p}_\sigma \leftarrow (1 - c_\sigma)\,\mathbf{p}_\sigma + \sqrt{c_\sigma(2-c_\sigma)\,\mu_{\mathrm{eff}}}\;\mathbf{C}^{-1/2}\frac{\mathbf{m}' - \mathbf{m}}{\sigma}
\]
gdzie $c_\sigma$ jest współczynnikiem kumulacji, a~$\mu_{\mathrm{eff}}$ efektywną wielkością populacji rodziców. Norma wektora $\mathbf{p}_\sigma$ jest porównywana z~oczekiwaną normą wektora losowego $\mathcal{N}(\mathbf{0}, \mathbf{I})$ --- jeśli jest większa, krok $\sigma$ rośnie; jeśli mniejsza --- maleje.
 
\textbf{Standardowa inicjalizacja} zakłada $\mathbf{p}_\sigma^{(0)} = \mathbf{0}$, co odpowiada brakowi wstępnej informacji o~kierunku przeszukiwania. Celem projektu jest zbadanie, jak \textbf{losowa inicjalizacja} tego wektora wpływa na zbieżność i~jakość wyników algorytmu.
 
\section{Sposób rozwiązania}
 
Projekt obejmuje implementację algorytmu CMA-ES w~dwóch wariantach:
\begin{itemize}
    \item \textbf{Wariant standardowy:} $\mathbf{p}_\sigma^{(0)} = \mathbf{0}$,
    \item \textbf{Wariant zmodyfikowany:} $\mathbf{p}_\sigma^{(0)} = \mathbf{v}$, gdzie $\mathbf{v}$ jest wektorem losowym.
\end{itemize}
Wektor losowy $\mathbf{v}$ będzie generowany na dwa sposoby, z~użyciem co najmniej dwóch różnych generatorów liczb pseudolosowych (PRNG), np.:
\begin{itemize}
    \item \texttt{numpy.random.default\_rng()} --- generator PCG64,
    \item \texttt{numpy.random.MT19937} --- generator Mersenne Twister.
\end{itemize}

\section{Planowane eksperymenty numeryczne}
 
Eksperymenty zostaną przeprowadzone na zbiorze klasycznych funkcji testowych o~różnej charakterystyce:
\begin{itemize}
    \item \textbf{Sphere} --- $f(\mathbf{x}) = \sum x_i^2$ --- wypukła, unimodalna, izotropowa,
    \item \textbf{Rosenbrock} --- wąska, zakrzywiona dolina, trudna dla metod gradientowych,
    \item \textbf{Rastrigin} --- silnie wielomodalna z~regularną siatką minimów lokalnych,
    \item \textbf{Ackley} --- wielomodalna z~jednym wyraźnym minimum globalnym.
\end{itemize}
 
Dla każdej konfiguracji (wariant inicjalizacji $\times$ generator $\times$ funkcja testowa) planowane jest wykonanie co najmniej 30 niezależnych uruchomień w~wymiarowościach $n \in \{2, 10, 30\}$. Mierzone metryki obejmą:
\begin{itemize}
    \item liczbę ewaluacji funkcji celu do osiągnięcia zadanej dokładności,
    \item końcową najlepszą wartość funkcji,
    \item krzywe zbieżności (wartość najlepsza vs.\ liczba ewaluacji).
\end{itemize}
 
Wyniki zostaną porównane z~użyciem statystyk opisowych (mediana, kwartyle) oraz testów istotności statystycznej (np.\ test Wilcoxona), aby ocenić, czy losowa inicjalizacja $\mathbf{p}_\sigma$ wpływa istotnie na działanie algorytmu.
 
\section{Wybór technologii}
 
\begin{itemize}
    \item \textbf{Język programowania:} Python 3.11+
    \item \textbf{Biblioteki obliczeniowe:} NumPy (operacje macierzowe, generatory PRNG), SciPy (ewentualnie algebra liniowa)
    \item \textbf{Wizualizacja:} Matplotlib, ewentualnie Seaborn
    \item \textbf{Analiza statystyczna:} SciPy (\texttt{scipy.stats} --- testy nieparametryczne)
    \item \textbf{Środowisko:} Jupyter Notebook (eksperymenty), skrypty \texttt{.py} (implementacja algorytmu)
    \item \textbf{Kontrola wersji:} Git
\end{itemize}
 
Implementacja CMA-ES zostanie napisana od podstaw w~celu pełnej kontroli nad inicjalizacją $\mathbf{p}_\sigma$ i~pozostałymi parametrami algorytmu. Jako punkt odniesienia posłuży biblioteka \texttt{cma} (autorstwa N.~Hansena), umożliwiająca walidację poprawności implementacji.
 

\end{document}
